\begin{frame}{Block B 개념 지도 읽는 법 — 박스와 축의 교차점}
  \begin{columns}[T]
    \begin{column}{0.53\textwidth}
      지도의 읽기는 8.3 Block A 지도와 같은 \textbf{두 기법}:
      \begin{itemize}
        \item \kb{(1) 질문은 박스가 아니라 박스와 축의 \emph{교차점}} —
          ``DQN의 리플레이 버퍼가 왜 PPO에는 없는가'' = DQN 박스와
          on/off-policy 축이 만나는 \textbf{엣지}; ``GAE의 $\lambda$를
          0.95로 올리면 어떻게 되나'' = PPO 박스와 편향-분산 축의
          엣지.
        \item \kb{(2) 엣지가 곧 ``왜''} — 박스 자체를 묻는 질문은
          드물다. 발표·면접에서 반복적으로 나오는 것은:
          \begin{itemize}
            \item \textbf{신경망으로 바꾼 이유}(Ch9, 표가 불가능해진
              상태공간)
            \item \textbf{클립을 하는 이유}(Ch11, 한 번의 업데이트가
              정책을 망가뜨리는 방어)
            \item \textbf{모방학습이 나빠지는 이유}(Ch12, 분포 이동의
              복합 오차)
            \item \textbf{sim-to-real이 필요한 이유}(Ch13, 현실 격차가
              \emph{편이}이기 때문)
          \end{itemize}
        \item \textbf{이 4개 엣지를 각각 한 문장}으로 말할 수 있으면,
          지도는 끝난 것이다.
      \end{itemize}
    \end{column}
    \begin{column}{0.45\textwidth}
      \centering
      \includegraphics[width=\linewidth]{ch16_3_rl_concept_map.png}
      \vskip0.3em
      {\footnotesize 중앙: ``정답 없이, 내가 만든 경험만으로, 로봇
      스케일까지 배우는가''. 위: 왜 신경망인가(Ch9$\sim$10). 가운데:
      어떻게 안정하게(Ch11$\sim$12). 아래: 어디에서 배우는가(Ch13
      $\sim$15). 가로지르는 축: 편향-분산($\lambda, \varepsilon$),
      탐험($\varepsilon$-greedy $\to$ UCT), on/off-policy, 모델
      유무. 오른쪽: 학기 끝(16.1$\sim$16.3) — 하나의 환경에서 정직한
      숫자로 증명.}
    \end{column}
  \end{columns}
\end{frame}
